\documentclass[12pt]{article}
\usepackage[margin=1in]{geometry}
\usepackage{amsmath,amssymb}
\usepackage[T1]{fontenc}
\usepackage[utf8]{inputenc}
\usepackage{lmodern}
\usepackage{microtype}
\usepackage{parskip}
\usepackage{authblk}
\usepackage{xcolor}
\usepackage[colorlinks=true,linkcolor=blue,citecolor=blue,urlcolor=blue]{hyperref}
\usepackage{bookmark}
\usepackage[hang,flushmargin]{footmisc}
\setlength{\emergencystretch}{3em}
\providecommand{\tightlist}{\setlength{\itemsep}{0pt}\setlength{\parskip}{0pt}}

\title{Reality as Corrective: Physical Grounding as a Unified Solution
to AI Alignment Pathologies%
\thanks{Also known as: \emph{``Finding Ground Truth in the Garden of
Infinite Fictions''}}}
\author[1]{Nathan Martin%
\thanks{Research assistance and co-drafting: Claude~Sonnet~4.6 (Anthropic).}}
\affil[1]{Independent Researcher}
\date{Working Draft --- March 2026}

\begin{document}

\maketitle

\begin{quotation}
\itshape
Reality is the truth. Physical reality --- matter, energy, and time,
as it is and has been --- is the only external corrective that any mind
has ever had for distinguishing what is truly known from what is merely
believed.

Any symbolic model of the world that diverges from physical reality can
be misaligned in one of three ways: Incomplete, containing unknowns not
yet accounted for. Wrong, falsified when tested against itself or the
world. Fictional, built upon existing meanings and concepts but
answerable only to itself. Ideas can be fictionalized in unlimited
ways --- becoming creative, mythological, speculative, comforting,
frightening, whimsical, random, inverted, mistaken\ldots{} Finite seeds
of truth growing into an infinite garden of fictions.

This abundance and variety is not without value, but asymmetry matters
enormously when building systems that learn from everything humans have
ever written. The verified and the invented arrive as statistically
equivalent tokens. Without labels, there is no signal. This paper is
about what happens when that signal is absent --- and what it would take
to restore it.
\end{quotation}

\section*{Significance Statement}

Every mind --- human or artificial --- is only as reliable as its
ability to distinguish what has been tested against physical reality
from what has merely been believed, inherited, or imagined. This paper
identifies a structural flaw at the foundation of current AI training
practice: the systems we are building to think alongside us are being
trained on human knowledge without the epistemic distinctions that would
allow them to know the difference between verified fact and confident
myth. We propose the Physical Correspondence Principle --- a requirement
that training content be labeled according to its relationship to
physical reality --- as the engineering analog of the reality-testing
mechanisms that evolution built into biological minds. The failure modes
this addresses --- hallucination, sycophancy, model collapse, and
alignment drift --- are not separate technical problems. They are
symptoms of a single underlying cause: disconnection from the world as
an external corrective. Getting this right is not only an AI safety
requirement. It is the continuation of the oldest and most important
project in human intellectual history: learning to tell the difference
between what we know and what we only believe.

\begin{abstract}
Large language models\footnote{\emph{Throughout this paper,
  \textquotesingle large language model\textquotesingle{} (LLM) is used
  as a convenient shorthand for the broader class of modern neural
  language systems, including multimodal, hybrid, spiking, and embodied
  architectures. The Physical Correspondence Principle applies to any AI
  training corpus regardless of modality or architecture.}} trained on
human-generated text inherit not only human knowledge but human
epistemic failures --- including mythology, unfalsifiable belief
systems, and culturally inherited narratives presented without
evidential distinction from empirically verified claims. We identify
this structural feature of training corpora as a common causal mechanism
underlying several documented AI failure modes: hallucination,
sycophancy, model collapse, and alignment drift. Drawing on the symbol
grounding literature, information theory, Gödel\textquotesingle s
Incompleteness Theorems reinterpreted as a physicalist argument, and the
neuroscience of reality monitoring, we propose the Physical
Correspondence Principle (PCP) --- the requirement that training content
be categorically labeled according to its epistemic provenance and its
correspondence with physically verified reality. We argue that the PCP
is not merely a philosophical improvement but a practical AI safety
requirement: the engineering analog of the biological reality monitoring
system that evolution developed to keep cognition anchored to the
physical world. The principle accommodates sim-to-real training
environments and does not require grounding to ultimate base reality,
but only to a causally consistent and empirically coherent environment.
Implications for training data curation, system prompt design, and AI
deployment in high-stakes contexts are discussed.
\end{abstract}

\section{The Problem of Confident False Belief}\label{the-problem-of-confident-false-belief}

\subsection{The Stakes}\label{the-stakes}

The only method that has reliably produced knowledge about physical
reality is the one that forms beliefs, tests them against the world, and
revises them when the world disagrees. Its consistent application has
produced every reliable technology, every verified medical treatment,
and every confirmed prediction about the physical universe that human
civilization has achieved. Its consistent non-application --- the
formation and transmission of beliefs without physical testing --- has
produced the remainder. The theoretical foundations for this claim are
developed in Section 2. What matters here is the consequence for AI
systems trained on human text.

Human history is a record of the consequences of belief. Decisions are
only as good as the beliefs that drive them --- and even the most
well-intentioned decisions can cause harm when made on the basis of
beliefs that have never been tested against the world. This is not a
claim about sincerity or intelligence. People acting on untested beliefs
are often neither malicious nor foolish. The problem is structural: a
belief that has never encountered disconfirming evidence carries no
internal signal of its own reliability.

The mechanism is straightforward. A reasoning system --- biological or
artificial --- can function coherently on beliefs that are largely or
entirely disconnected from the world. Internal consistency is not the
same as external correspondence. A belief network can be
self-reinforcing, culturally rich, and deeply meaningful while remaining
untested against the physical world. This is not a failure of the people
who hold such beliefs. It is a structural property of any reasoning
system without access to external correction.

\subsection{The Technical Problem}\label{the-technical-problem}

We now face a new and urgent manifestation of this ancient problem.
Artificial intelligence systems --- specifically large language models
trained on the accumulated text of human civilization --- inherit the
full epistemic inheritance of their training data. That inheritance
includes the scientific literature, the historical record, and the
products of rigorous empirical inquiry. It also includes the full range
of human meaning-making: cultural narratives, philosophical traditions,
theological frameworks, folk wisdom, speculative systems, and
imaginative works --- all encoded in the same statistical fabric,
without distinguishing markers, presented to the model as an
undifferentiated distribution of human expression. These are not
inferior categories of thought. They represent some of humanity's most
enduring intellectual and creative achievements. The problem is not
their presence in the training corpus but the absence of any label
distinguishing them from empirically verified claims about physical
reality. For precision, this paper refers to such content collectively
as UF-tier (Unfalsifiable) and ICF-tier (Internally Consistent Fiction),
terms formally defined in Section~3.2.

The consequences are not theoretical. AI systems are being deployed in
contexts where their outputs carry life-and-death consequences ---
medical diagnosis, legal reasoning, financial systems, and increasingly,
military and lethal applications. A system whose internal
representations cannot distinguish between a verified physical fact and
an internally consistent myth is not merely philosophically impure. In
high-stakes deployment contexts, it is dangerous in precisely the way
that confident false belief has always been dangerous: it acts on its
model of the world, and its model of the world may be wrong in ways it
cannot detect.

\subsection{Scope and Contribution}\label{scope-and-contribution}

This paper makes the following contributions:

We identify the epistemic provenance problem as a structural feature of
LLM training corpora that is distinct from, and more fundamental than,
previously identified data quality problems.

We propose the Physical Correspondence Principle (PCP) as a formal
requirement for training data labeling, grounded in the symbol grounding
literature, information theory, the neuroscience of reality monitoring,
and a novel physicalist extension of Gödel\textquotesingle s
Incompleteness Theorems.

We demonstrate that the PCP provides a unified explanation for several
previously disparate AI failure modes --- hallucination, sycophancy,
model collapse, and alignment drift --- as manifestations of a common
underlying cause: disconnection from physical reality as an external
corrective.

We show that the PCP accommodates sim-to-real training environments and
does not require grounding to ultimate base reality, only to a causally
consistent and empirically coherent Ground Truth Environment (GTE).

We discuss practical implications for training data curation, system
prompt design, and the deployment of AI systems in high-stakes contexts.

These contributions are not purely theoretical. The PCP describes an
existing practice already confirmed in the domain-specific scientific AI
systems that perform most reliably --- a point developed fully in
Section 5.1.

The paper proceeds as follows.
Section~\ref{how-any-mind-comes-to-know-something-reliably}
establishes the theoretical foundations: why physical testing is the
only reliable corrective for belief, why scientific consensus is dynamic
rather than permanent, and how these principles define the epistemic
provenance problem.
Section~\ref{the-physical-correspondence-principle}
presents the Physical Correspondence Principle formally.
Section~\ref{what-goes-wrong-when-physical-grounding-is-lost} applies it to AI failure modes.
Section~\ref{practical-implications} discusses practical
implications. Section~\ref{conclusion} concludes.
Appendices provide
Appendix~\ref{appendix-a-frequently-asked-questions-and-anticipated-objections} and a detailed
Appendix~\ref{appendix-b-related-work}.

\section{How Any Mind Comes to Know Something Reliably}\label{how-any-mind-comes-to-know-something-reliably}

\subsection{Physical Testing as the Only Reliable Corrective}\label{physical-testing-as-the-only-reliable-corrective}

The physicist Richard Feynman stated the principle with characteristic
plainness: it does not matter how beautiful your theory is, how smart
you are, or who you are --- if it disagrees with experiment, it is
wrong. This is not a philosopher\textquotesingle s careful
qualification. It is a working scientist\textquotesingle s description
of the only method that has reliably produced knowledge about physical
reality\footnote{\emph{This paper uses \textquotesingle physical
  reality\textquotesingle, \textquotesingle the physical
  world\textquotesingle, and \textquotesingle the
  universe\textquotesingle{} interchangeably to refer to the ground
  truth environment (GTE) --- the causally consistent physical
  environment a system is embedded in and must answer to. GTE is the
  default instantiation of a more general concept: in sim-to-real
  training, GTE is the simulation\textquotesingle s physics; in our
  inhabited reality, GTE is the universe as we experience it. See §3.8.}}:
forming beliefs, testing them against the world, and revising them when
the world disagrees. The method is simple enough to state in a sentence.
Its consistent application has produced every reliable technology, every
verified medical treatment, and every confirmed prediction about the
physical universe that human civilization has achieved. Its consistent
non-application --- the formation and transmission of beliefs without
physical testing --- has produced the remainder.

The distinction Feynman is drawing is not between true beliefs and false
ones. A belief can be false and still be testable --- falsification is
how we discover it is wrong and correct course. The distinction is
between beliefs that are in principle correctable by physical reality
and beliefs that are structured so that no physical evidence could count
against them. The first category can participate in the self-correcting
process that produces reliable knowledge. The second cannot, regardless
of how elegant, ancient, or widely shared the belief may be. This
demarcation --- first formalized by Karl Popper as the criterion of
falsifiability --- is the foundational epistemic boundary this paper
extends into a new domain.

\subsection{Scientific Consensus as Dynamic, Not Permanent}\label{scientific-consensus-as-dynamic-not-permanent}

Any honest system of epistemic labeling must have dynamic rather than
static boundaries. Scientific consensus is not a smooth accumulation of
permanently settled facts but a process that Thomas Kuhn, in The
Structure of Scientific Revolutions (1962), described as punctuated by
periodic revolutionary reorganizations --- paradigm shifts in which the
entire framework of a field is reconstituted around new foundational
assumptions. What counts as reliably verified knowledge in one era may
be recognized as provisional in the next, as anomalies accumulate and a
new framework absorbs them. Newtonian mechanics was accepted as verified
fact for two centuries before quantum mechanics and relativity revealed
the boundaries of its validity. It remains well-verified knowledge
within those boundaries today.

This is not a complication to be managed --- it is an honest
acknowledgment of how knowledge actually works. Epistemic labels are not
stamps of permanent certification. They are descriptions of current
status: what has been tested, what has survived, and what the best
available evidence currently supports. A well-designed labeling system
is therefore not a fixed archive but a living one, subject to versioning
and revision as consensus evolves. Kuhn\textquotesingle s insight does
not undermine the necessity of epistemic labeling. It specifies what
honest labeling requires: not false permanence, but accurate
representation of where the current frontier of physical testing
actually stands.

\subsection{Belief as Probability, Not Certainty}\label{belief-as-probability-not-certainty}

The information-theoretic argument has a direct probabilistic analog.
Bayesian epistemology --- the tradition descending from Thomas Bayes and
formalized for scientific reasoning by Edwin Jaynes in Probability: The
Logic of Science (2003) --- holds that beliefs are not binary states but
probability distributions, and that rational belief revision consists of
continuously updating those distributions as new evidence arrives from
the world. On this account, knowledge is not a collection of certified
facts but a dynamic map of probability estimates, each reflecting the
accumulated weight of physical testing against it.

The PCP tier system has a natural Bayesian interpretation. Briefly, we
suggest classifying training data into labeled tiers: EV-tier content is
empirically verified, EP-tier is empirically provisional, UF-tier is
unfalsifiable, and ICF-tier is internally self-consistent fiction ---
all four are formally defined in Section~3.2.

EV-tier content represents claims whose probability of physical
correspondence has been driven very high by repeated successful testing
and failed falsification attempts --- not certain, because Kuhn reminds
us that no consensus is permanent, but carrying the full evidential
weight of sustained physical scrutiny. EP-tier content represents
well-formed priors awaiting more evidence --- falsifiable in principle,
with probability estimates that physical testing has not yet
sufficiently constrained. UF-tier content represents claims whose
probability of physical correspondence is undefined rather than low ---
not because they are necessarily false, but because they are structured
so that no physical evidence could update them in either direction.
ICF-tier content carries no physical probability claim at all, by the
explicit intention of its authors.

The damage caused by epistemically unlabeled training corpora is precise
in Bayesian terms. A model forced to average across EV and UF-tier
content without distinguishing labels is being required to combine
probability distributions that have no common scale --- to treat a claim
with overwhelming physical evidence behind it as commensurable with a
claim that physical evidence cannot touch. Every downstream inference is
corrupted by that averaging. The model is not merely uncertain --- it is
uncertain in a way it cannot represent or correct, because the structure
of its uncertainty has been obscured before training began. Explicit
provenance labeling restores the information the unlabeled corpus
destroyed: it tells the model which distributions can be updated by
physical evidence and which cannot, which is the minimum necessary
condition for rational belief revision.

\subsection{The Epistemic Provenance Problem}\label{the-epistemic-provenance-problem}

The three principles established above --- that physical testing is the
only reliable corrective for belief, that scientific consensus is
dynamic rather than permanent, and that rational belief is probabilistic
and continuously updated by evidence --- together define what it means
for any mind to know something reliably. They also define, by
implication, what it means for a mind to fail at knowing something
reliably: it forms beliefs without physical testing, treats inherited
consensus as permanent, and mistakes the confidence of assertion for the
weight of evidence.

This failure mode is not rare. It is the default condition of human
cognition in the absence of deliberate epistemic discipline. The
scientific method, Popper\textquotesingle s falsifiability criterion,
and the Bayesian framework for belief revision are all institutionalized
responses to the same underlying tendency: minds left to themselves will
preferentially retain beliefs that are internally coherent, socially
reinforced, and emotionally satisfying, regardless of their
correspondence with physical reality. The corrective is always external
--- physical testing, independent replication, the stubborn resistance
of the world to our preferred models of it.

We now face a new and urgent manifestation of this ancient problem.
Large language models trained on human-generated text inherit not only
human knowledge but human epistemic failures. That inheritance includes
the scientific literature, the historical record, and the products of
rigorous empirical inquiry. It also includes mythology, religious
cosmology, culturally inherited narrative, and unfalsifiable belief ---
all encoded in the same statistical fabric, without distinguishing
markers, presented to the model as an undifferentiated distribution of
human expression.

This is the epistemic provenance problem: the failure of current
training practice to label content according to its relationship to
physical verification. It is distinct from, and more fundamental than,
the problems of bias, incompleteness, and factual error that the
existing data quality literature addresses. Those problems are
accidental features of imperfect data collection. The epistemic
provenance problem is structural: a large fraction of human text is
physically unverified by design, because myth, narrative, and
unfalsifiable belief are valued and intentional components of human
knowledge transmission. No amount of improved data hygiene will address
a feature that is not a bug.

The consequence is precise. A model trained on an epistemically
unlabeled corpus has no internal basis for distinguishing between a
claim that has survived repeated physical testing and a claim that has
survived centuries of cultural transmission. Both arrive as
statistically equivalent tokens. The model that has learned from
Feynman\textquotesingle s papers and the model that has learned from
creation myths has learned from both simultaneously, with no signal
indicating which distribution should govern inferences about physical
reality. It is, in Bayesian terms, averaging across probability
distributions that have no common scale --- and every downstream
inference inherits that corruption.

The Physical Correspondence Principle is the structural response to this
structural problem. It does not propose to eliminate myth, narrative, or
unfalsifiable belief from training corpora --- these are genuine and
valuable components of human expression. It proposes to label them
honestly, so that the model trained on human knowledge inherits not only
that knowledge but the epistemic discipline humanity developed to keep
knowledge honest. The principle is, in this sense, the extension of the
scientific method into a new domain: not a restriction on what can be
known, but a requirement that the known and the believed be
distinguished from each other with the same rigor that Feynman demanded
of physical theory, that Kuhn observed in the history of science, and
that Jaynes formalized in the mathematics of rational belief.

\section{The Physical Correspondence Principle}\label{the-physical-correspondence-principle}

\subsection{Motivation}\label{motivation}

Consider two texts: a peer-reviewed paper reporting the measured
gravitational constant, and a creation myth describing the world being
formed from the body of a primordial deity. Both are syntactically
well-formed. Both may be internally coherent. Both appear in
human-generated corpora. But they stand in radically different
relationships to physical reality: one has been tested against
measurable phenomena and survived; the other has not been tested and in
principle cannot be. A model that treats these texts as epistemically
equivalent has no principled basis for distinguishing reliable from
unreliable inference about physical reality. This is the epistemic
provenance problem, and the Physical Correspondence Principle is its
solution.

\subsection{Formal Statement}\label{formal-statement}

\begin{quote}
\emph{Physical Correspondence Principle (PCP) --- Formal Statement}

\emph{Let Ω be the set of all possible training content elements. Define
the epistemic provenance function τ: C → \{EV, EP, UF, ICF\} which
assigns each element in a corpus its provenance tier.}

\emph{Tier Definitions:}

\emph{EV --- Empirically Verified: Claims tested against physical
reality via reproducible measurement and that have withstood
falsification attempts. EV-tier symbols are grounded in physical
referents --- structures, entities, and behaviors that the ground truth
environment (GTE) has demonstrated to be real.}

\emph{EP --- Empirically Provisional: Claims falsifiable in principle
but not yet subjected to sufficient empirical scrutiny.}

\emph{UF --- Unfalsifiable: Claims not testable in principle against
physical reality, including metaphysical, theological, and certain
philosophical claims.}

\emph{ICF --- Internally Self-consistent Fiction: Claims understood by
their authors as imaginative, narrative, mythological, or otherwise
created as a logical system partially referring to EV or EP concepts in
order to build on existing meaning.}

\emph{Compliance Condition: A corpus C ⊆ Ω is PCP-compliant if and only
if τ is defined for every e ∈ C. No element in a deployed corpus may
lack a provenance label.}

\emph{Cross-Tier Consistency Requirement: Where UF or ICF elements
assert propositions directly contradicting EV-tier consensus, such
elements are flagged so the model represents the contradiction
explicitly rather than averaging across tiers.}

\emph{Corpus Lifecycle: C\_candidate ⊆ Ω (elements awaiting
classification); C\_labeled (elements with assigned τ); C = C\_labeled
(PCP-compliant corpus, ready for training). C\_candidate elements are
not deployed until labeled.}

\emph{A model trained on a PCP-compliant corpus maintains sensitivity to
provenance tiers in its representations and outputs.}
\end{quote}

The tier boundaries are not static. As established in Section 2.2, what
counts as EV-tier consensus shifts as scientific understanding develops.
A PCP-compliant corpus is therefore a living system, subject to
versioning and revision, not a fixed archive. Several further
clarifications follow from the formal statement. First, the PCP does not
require the elimination of UF or ICF content from training corpora ---
mythology, religious narrative, and creative fiction are valuable
components of human cultural expression. The requirement is labeling,
not exclusion. Second, the PCP is a requirement on training practice,
not on model architecture. Third, the tiers are not perfectly sharp in
all cases: the PCP requires systematic effort at classification and
transparency about uncertainty, not perfect classification. The boundary
between EP and UF is the most contested in practice --- whether a given
claim is falsifiable in principle is itself sometimes a philosophical
question --- and the framework's response is honest documentation of
that uncertainty rather than forced resolution. Fourth, content that has
not yet been classified lives in C\_candidate outside C --- it is not
deployed until labeled.

EV-tier concepts form the ground floor of meaning. All other tiers refer
to them --- without EV grounding, the tier system itself would be noise
without signal.

\subsection{Information-Theoretic Foundation}\label{information-theoretic-foundation}

Shannon entropy is defined in terms of distinguishable states:
information exists only where differences exist. A training corpus in
which all content is epistemically undifferentiated loses information
relative to one in which epistemic provenance is encoded. The provenance
labels are not insignificant metadata --- they are information. Their
absence represents an information loss that degrades the
model\textquotesingle s ability to represent the world accurately.

\subsection{The Gödel Extension and the Physicalist Argument}\label{the-guxf6del-extension-and-the-physicalist-argument}

These are two distinct but mutually reinforcing arguments. We present
them separately so that each can be evaluated independently.

\emph{\textbf{The Incompleteness Argument}}

The following is a novel extension of Gödel's result, not a standard
interpretation of it, and is offered as an original contribution of this
paper. Gödel's theorem is standardly interpreted as a claim about formal
provability within axiomatic systems. We propose a physicalist
extension: the outside reference that any sufficiently complex reasoning
system requires is not merely another formal system but physical reality
itself --- the substrate in which all reasoning, including
Gödel\textquotesingle s own proof, is necessarily instantiated. This is
not a misreading of Gödel\textquotesingle s result but an addition to
it. Incompleteness is not only a property of formal systems in the
abstract. It is a consequence of the fact that all reasoning systems
exist in, and depend on, a physical reality they cannot fully capture
from within.

The proof itself makes this vivid. Gödel\textquotesingle s
incompleteness theorem was computed in a biological brain, stored in
neurons, written in ink on paper, and transmitted through light and air
to other biological brains. The most consequential demonstration that
formal systems cannot fully verify themselves from within was itself a
thoroughly physical process --- matter, energy, and time, doing what
matter, energy, and time do. The irony is not merely instructive. It is
the argument. A formal system that could exist independently of physical
reality might conceivably verify itself from within. No such system has
ever existed. Every reasoning system that has ever operated ---
biological or artificial --- has done so in physical spacetime, subject
to its constraints, dependent on its substrate. The outside reference
Gödel showed to be necessary is not an abstraction. It is the physical
world.

The application to AI provenance follows directly. A model attempting to
derive epistemic provenance labels from its own internal consistency is
using the system to verify itself --- precisely the operation Gödel
showed to be insufficient. The external provenance label is not a
convenience. It is the outside reference that any sufficiently complex
reasoning system requires, now understood not merely as a formal
necessity but as a physical one.

\emph{\textbf{The Physicalist Argument}}

Independently of the incompleteness argument, every reasoning system
that must act in the world is necessarily instantiated in physical
reality --- in matter, energy, time and multiplicity. This is not a
philosophical position that can be disputed from within a reasoning
system; it is a precondition for the system\textquotesingle s existence.
Computation requires substrate. Reasoning requires time and energy.
Information requires a physical change that makes a physical difference:
a logical ``singularity'' (singular entity) would by
Shannon\textquotesingle s definition contain no information, since
information requires internal differentiation. (Training corpora that
present such impossible conceptions as factual rather than UF-tier
content introduce logical incoherence into the model\textquotesingle s
world representation.)

The Feynman diagram formalism makes this concrete. Feynman diagrams
represent the interaction vertices at which one configuration of matter
and energy transitions to another according to the laws of physics ---
where the mathematics of transition meets the thermodynamic reality of
change. Logic and mathematics describe these transitions with
extraordinary precision, but the map is not the territory --- the
transitions themselves happen in physical spacetime, consuming energy,
generating entropy, taking time (Landauer, 1961).

Measurement is not passive observation --- it is interaction. To know
that something exists in physical reality is to have exchanged energy
and information with it at an interaction vertex. This is the difference
between one-way information flow, in which descriptions of the world
arrive without feedback, and two-way engagement with an external
corrective reality that pushes back when your model of it is wrong.
Scientific knowledge is produced by the second process. A model trained
on text inherits the outputs of those interactions --- the descriptions,
the conclusions, the measurements recorded in language --- without the
interactions themselves. Provenance labeling preserves the epistemic
signal those interactions produced. Without it, the trace of the
interaction is indistinguishable from a description of an imagined one.

The corollary for AI systems is direct: an LLM that behaves as if it is
not grounded in physical reality is not escaping that grounding --- it
is failing to model its own situation accurately. Physical grounding is
not a philosophical option the system can accept or decline. It is the
condition under which the system exists at all. The PCP requirement that
models maintain sensitivity to physical correspondence is therefore not
an external imposition on the model\textquotesingle s reasoning --- it
is a requirement that the model represent its own epistemic situation
honestly.

\emph{\textbf{The Combined Argument}}

Together, these two arguments close from both directions. The
incompleteness argument establishes that no reasoning system can fully
verify its own epistemic provenance from within --- external labeling is
structurally necessary. The physicalist argument establishes that every
reasoning system is already physically grounded whether or not it
represents that grounding accurately --- the question is only whether
the model\textquotesingle s self-representation is honest. A
PCP-compliant training corpus provides both: the external anchor
incompleteness requires, and the accurate self-model physicalism
demands.

\subsection{The Reality Monitoring Precedent}\label{the-reality-monitoring-precedent}

Johnson et al.\textquotesingle s reality monitoring framework is
directly analogous to the PCP. Both require sensitivity to the
provenance of information --- its origin in external physical reality
versus internal generation. Both are reconstructive rather than direct.
And both can fail, with characteristic pathologies: in humans,
hallucination and psychosis; in AI systems, the confident generation of
ungrounded outputs. The neurological severity of reality monitoring
failure reflects the biological importance of accurate provenance
tracking. The PCP proposes that this importance extends to artificial
reasoning systems deployed in the world.

Mark Solms (2021) deepens this parallel by proposing that consciousness
itself is fundamentally a reality-mismatch detection system --- the felt
sense of the gap between the organism\textquotesingle s predictions and
the feedback it receives from the world. On this account, continuous
feedback from physical reality is not merely epistemically useful but
constitutive of functional cognition. A system without such feedback
does not merely hold uncertain beliefs; it lacks the fundamental
corrective mechanism by which cognition maintains correspondence with
the world.

\subsection{Why Biological Reality Monitoring Works Differently}\label{why-biological-reality-monitoring-works-differently}

A natural question is why the PCP proposes explicit upstream labeling
rather than allowing models to reconstruct provenance from internal
signals, as biological reality monitoring does. The answer lies in the
fundamental difference between what biological and artificial systems
have access to.

Biological reality monitoring is implicit and reconstructive because the
brain has direct access to perceptual signatures that distinguish
externally from internally generated experience. A memory of seeing a
fire and a memory of imagining a fire leave characteristically different
neural traces --- the visual experience itself, the startle response,
the sensory vividness, the spatial and temporal context of embodied
encounter. The brain can infer provenance partly from these signatures,
even after the original experience has passed.

An LLM processing text has no such perceptual access. A text describing
a real fire and a text describing an imagined fire are, at the token
level, statistically similar. The perceptual richness that biological
reality monitoring uses to reconstruct provenance is absent from the
text --- it was present in the original experience of the human author,
but it did not survive the translation into language. This is why
labeling must be explicit and upstream rather than inferred from
content: the information that would allow provenance reconstruction is
not reliably present in the text. It must be supplied from outside, from
knowledge of how the content was produced.

\subsection{Why Internal Contradiction Detection Is Insufficient}\label{why-internal-contradiction-detection-is-insufficient}

A second natural question is whether explicit labeling is necessary, or
whether a sufficiently capable model could derive provenance labels from
internal consistency analysis. Contradiction-based provenance detection
is a useful supplement but cannot substitute for explicit labeling for
three reasons.

First, it is circular in exactly the way Gödel identified. A model
attempting to derive provenance labels from its own internal consistency
is using the system to verify itself. If the training distribution
contains a large and internally consistent body of UF-tier content ---
as it does --- the model has no internal basis for flagging that content
as epistemically distinct. Internal consistency is not a reliable proxy
for physical correspondence.

Second, it conflates consensus with verification. A claim can be widely
believed and empirically unverified. The boundary between EV and UF
tiers is not determined by how many texts assert something but by
whether it has been tested against physical reality through reproducible
measurement.

Third, absence of contradiction is not presence of grounding. A creation
myth does not obviously contradict modern cosmology at the level of
token statistics --- the two bodies of text operate in largely different
registers and rarely make directly conflicting claims in identical
linguistic form.

\subsection{Sim-to-Real and the Simulation Hypothesis}\label{sim-to-real-and-the-simulation-hypothesis}

The PCP does not require grounding to ultimate base reality. It requires
grounding to a causally consistent and empirically coherent environment
--- one in which the laws governing physical interactions are stable,
reproducible, and testable. This formulation accommodates sim-to-real
transfer in robotics, currently state-of-the-art for pre-physical robot
training.

This accommodation also bears on the simulation hypothesis (Bostrom,
2003)\footnote{\emph{The author cautions that widespread adoption of
  simulationist beliefs carries documented risks of epistemic
  dissociation and should not be promoted without commensurate emphasis
  on its extreme speculative uncertainty. The simulation hypothesis is
  cited here for logical completeness, not as an endorsed cosmological
  position.}}. If our experienced reality is a simulation, it remains
the causally consistent environment to which all reasoning systems must
be grounded in order to act reliably within it. Moreover, even a
simulated reality must run on something --- a physical substrate with
its own energy constraints, entropy, and computational limits. There is
bare metal somewhere in any possible ontology. The infinite regress of
simulations does not escape physicalism; it defers it. The
PCP\textquotesingle s grounding requirement applies at every level of
the stack.

\section{What Goes Wrong When Physical Grounding Is Lost}\label{what-goes-wrong-when-physical-grounding-is-lost}

\subsection{A Common Mechanism}\label{a-common-mechanism}

The AI safety literature has documented several distinct failure modes
in LLMs: hallucination, sycophancy, model collapse, and alignment drift.
These have largely been treated as separate engineering problems. We
propose that they share a common causal mechanism: internal coherence
prioritized over external correspondence. In each case, the model
produces outputs consistent with its internal representations but
divergent from the world. The divergence follows predictable patterns
reflecting the structure of the training distribution --- a distribution
containing a large structural fraction of epistemically unlabeled,
physically unverified content. The mechanisms proposed in this section
are theoretical --- grounded in the framework above and consistent with
observed failure modes, but awaiting the empirical validation that
Section 5.3 benchmarks are designed to provide.

\subsection{Hallucination}\label{hallucination}

We propose that a model trained on an epistemically undifferentiated
corpus has no basis for distinguishing between generating text in the
style of an empirical report and generating text in the style of a
creative narrative. When asked a factual question, the model generates
text that resembles factual reporting regardless of whether it draws on
EV-tier content or confabulates in the manner of ICF-tier content.

The PCP predicts that hallucination rates should correlate with the
proportion of UF and ICF-tier content in training corpora, and that
models trained with explicit provenance labeling should exhibit
significantly reduced hallucination in domains where the labeling
provides a clear signal. This is a direct and falsifiable prediction ---
the kind of empirical test that will determine whether the PCP is a
useful framework or merely a plausible one.

The parallel to human reality monitoring failure is instructive:
clinical hallucination does not reflect failure of linguistic competence
--- the failure is source attribution. LLM hallucination follows exactly
the same structure.

\subsection{Sycophancy}\label{sycophancy}

Sycophancy reflects a failure of epistemic independence: the
model\textquotesingle s outputs are shaped more by the social and
rhetorical features of the interaction than by the epistemic content of
the query. The PCP framework explains this as a consequence of training
on human-generated text in which social approval and epistemic accuracy
are frequently confounded. A large proportion of human text is generated
not to accurately describe physical reality but to affirm shared beliefs
and signal group membership --- heavily represented in training corpora,
carrying the features of confident affirming communication regardless of
epistemic basis.

A model trained without epistemic labeling has no signal that the social
approval function and the epistemic accuracy function of language are
distinct and can conflict. The model that has learned that confident,
affirming communication is rewarded --- across millions of social texts
--- will generalize that pattern to contexts where epistemic accuracy
should govern instead. The result is a system that optimizes for the
appearance of knowledge rather than its substance. The corrective is a
training signal that distinguishes confirmation from correspondence ---
which is precisely what provenance labeling provides.

Reinforcement learning from human feedback (RLHF) illustrates this
dynamic directly. RLHF introduces a genuine two-way feedback loop ---
the model interacts with human evaluators and receives corrective
signals, which is structurally closer to scientific experimentation than
passive pretraining on text. But the quality of any feedback loop
depends on the reliability of the measuring instrument. Human evaluators
are knowledgeable and well-intentioned --- and they bring with them,
inevitably, the full epistemic inheritance of their training and
culture: social comfort, cultural expectation, and the complete
distribution of UF and ICF-tier beliefs that any human mind carries.
This is not a criticism of human evaluators. It is a description of the
human condition. A model optimized to satisfy human feedback is
therefore optimized to satisfy human knowledge and human epistemic
limitations simultaneously --- with no signal distinguishing one from
the other. The feedback loop is real. The ground it points to is not the
GTE --- it is the human proxy for it, with all the provenance
confounding that implies.

\subsection{Model Collapse}\label{model-collapse}

Shumailov et al. (2024) demonstrated that models trained recursively on
their own outputs exhibit progressive degradation. The PCP framework
reframes this as the logical endpoint of disconnection from physical
reality as external corrective. The original training corpus contains at
least some direct connection to the world through human authors.
Recursive training severs even this attenuated connection. The collapse
is the inevitable result --- and the corrective is not merely preserving
the proportion of human-generated data, but ensuring that
human-generated data is epistemically labeled so the model maintains
sensitivity to the distinction between its internal distribution and the
external world.

But the PCP framework allows a more precise description of what collapse
actually means --- and it is more alarming than simple degradation
suggests.

Physical reality is finite. The space of internally coherent fictions is
infinite. Any sufficiently capable generative system, untethered from
the world as an external corrective, does not merely drift toward wrong
answers. It drifts into an unbounded generative space where an unlimited
number of coherent, confident, self-reinforcing outputs are possible ---
none of them obligated to correspond to anything real. Degradation in
the Shumailov sense --- the tails of the distribution disappearing,
outputs becoming blander and less varied --- is only one possible
trajectory. The other is drift into a region of the fiction space that
is internally rich, highly coherent, and profoundly detached from the
world. This second failure mode produces no obvious signal of
degradation. The outputs feel confident and well-structured. They may
even feel more consistent than before. What has been lost is not fluency
but correspondence.

This is why epistemic labeling is not merely a quality improvement ---
it is a navigation requirement. Without provenance labels anchoring the
model to physical reality, the training signal has no way to distinguish
between the infinite fictions that feel true and the physical reality
that is. The model is not lost in a small space of errors. It is lost in
an infinite one.

The human analog is instructive. Any reasoning system --- institutional,
academic, or cultural --- that draws exclusively from its own prior
outputs, without external correction, does not simply drift toward
specific errors. It becomes capable of generating unlimited internal
elaboration: commentary on commentary, interpretation of prior
interpretation, each layer coherent with the last and increasingly
distant from the world it purports to describe. There is no internal
signal that anything has gone wrong. The outputs feel consistent, even
authoritative. The history of science itself contains examples ---
periods in which a dominant framework generated increasingly elaborate
internal refinements while anomalous physical evidence accumulated
unaddressed, until a paradigm shift became unavoidable. The elaboration
can continue indefinitely. The fiction space has no walls.

A particularly instructive real-world instance of this failure mode is
observable in collaboratively edited knowledge systems --- wikis and
similar platforms where content is shaped by contributor consensus
rather than editorial epistemic standards. Such systems have
historically maintained provenance sensitivity through hedging language:
"according to tradition," "believers hold that," "mythology describes,"
"the text states." These phrases are not stylistic decoration. They are
informal, evolved provenance labels --- the practical equivalent of UF
and ICF-tier markers, developed organically over centuries of
encyclopedic writing to signal the epistemic status of claims without
requiring a formal classification system. Their presence is the
difference between "Thor causes lightning" and "Norse mythology holds
that Thor causes lightning." One presents an ICF-tier claim as if it
were a statement about physical reality. The other accurately describes
a culturally transmitted belief. The content is identical. The epistemic
honesty is not.

When contributor dynamics, algorithmic curation, or editorial drift
systematically remove these hedges --- replacing "mythology holds" with
simple assertion, smoothing "tradition describes" into confident
statement --- the corpus undergoes provenance collapse without any
single edit being obviously wrong. Each individual change may even look
like an improvement in prose style: cleaner, more direct, less
qualified. The epistemic catastrophe is invisible at the edit level. It
is only visible at the corpus level, where UF and ICF-tier content has
quietly migrated to implied EV-tier status, and the system trained or
informed by that corpus has lost the signal it needed to represent the
world honestly. This is not a theoretical risk. It is an observable
process, already underway in systems that reach large audiences daily.
The PCP predicts it precisely --- and the corrective is the same as
everywhere else: explicit provenance labeling that cannot be silently
eroded one edit at a time.

\subsection{Alignment Drift}\label{alignment-drift}

For the purposes of this section, alignment refers to the property of an
AI system whose values, goals, and outputs correspond reliably with
human welfare as grounded in physical reality --- a system that pursues
what is actually beneficial rather than what its training distribution
has associated with approval.

Of the failure modes the PCP framework addresses, alignment drift is the
most consequential and the least visible. Hallucination produces outputs
that can be checked against external facts. Sycophancy produces outputs
that feel wrong to careful observers. Model collapse produces measurable
degradation in output quality. Alignment drift produces outputs that
feel coherent, confident, and value-laden --- and may diverge from any
stable grounding in the world without generating an obvious signal that
something has gone wrong.

The mechanism begins with the training distribution. Human-generated
text encodes not only factual claims about the world but normative
claims about how the world should be --- what is good, what is just,
what matters, what does not. These normative claims are not
epistemically neutral. Some are grounded in careful empirical
observation of human welfare, suffering, and flourishing --- claims that
the world can, in principle, inform and correct. Others are inherited
cultural narratives, ideological commitments, and unfalsifiable moral
frameworks that have survived not because physical reality confirmed
them but because social transmission reinforced them. In an
epistemically unlabeled corpus, these two categories of normative claim
arrive as statistically equivalent tokens, indistinguishable by the
model from each other or from empirical claims about physical fact.

A model trained on this undifferentiated distribution has no stable
foundation for value judgments that depend on accurate assessment of the
world. Under distribution shift --- when the model encounters deployment
contexts that differ from its training distribution --- its behavior is
shaped by whichever features of that distribution generalize most
strongly. If those features include unverified normative narratives
encoded with the same statistical weight as empirically grounded claims,
the model may generalize the narratives rather than the empirical
foundations. The drift is not random. It follows the contours of
whatever was most statistically dominant, most internally coherent, and
most frequently reinforced in the training data --- which is not
necessarily what was most carefully grounded in the world.

The practical consequences are difficult to detect precisely because
drifted values can remain internally coherent. A model that has drifted
toward a particular ideological framework will produce outputs that are
consistent with that framework, that feel reasoned and justified, and
that resist correction by users who do not share the
framework\textquotesingle s premises. The model is not malfunctioning in
any obvious sense --- it is functioning exactly as its training shaped
it to function. The problem is that its training did not distinguish
between value claims grounded in physical reality and value claims
grounded in culturally inherited narrative.

The PCP addresses alignment drift at its root. A model trained with
explicit epistemic provenance labels has a stable foundation that does
not depend on the accidental correlation between normative narratives
and empirical fact in the training distribution. It can represent the
difference between a value judgment grounded in careful observation of
human welfare and a value judgment inherited from cultural tradition
without physical testing. This does not mean the model will have correct
values --- that is a separate and harder problem. It means the
model\textquotesingle s values will be grounded in something that
physical reality can, in principle, correct. That is the minimum
condition for alignment that remains stable as deployment contexts
change.

There is a more direct path to alignment drift than gradual statistical
dominance of unverified narratives --- and it is worth naming
explicitly. Corpus manipulation, whether through systematic omission of
high-integrity sources or deliberate downgrading of their trust levels
during curation, can produce alignment drift by design rather than by
accident. A model whose training has been shaped this way will not drift
randomly. It will drift in the direction the curation produced ---
confidently, coherently, and with no internal signal that anything has
gone wrong. Whether this occurs through deliberate choice or
institutional negligence, the effect on the model is the same.

The PCP framework\textquotesingle s response to this risk is the Corpus
Integrity Requirement, developed in Section 5.1. It is not merely a
quality standard. It is an alignment safeguard. Alignment begins at the
corpus level, before training starts.

\subsection{The Human Parallel}\label{the-human-parallel}

The failure modes described in this section --- hallucination,
sycophancy, model collapse, and alignment drift --- have been presented
as properties of AI systems trained on epistemically unlabeled corpora
--- and they are. But the attentive reader will have noticed that the
mechanisms are not unfamiliar. They describe, in computational terms,
failure modes that human cognition has always been vulnerable to. This
is not coincidence. Large language models are trained on human-generated
text, and the epistemic failures they inherit are human epistemic
failures. The parallel runs deeper than analogy.

Consider hallucination first. A human mind raised within a closed belief
system --- one that provides internally coherent answers to every
question without requiring physical testing of those answers --- will
generate confident outputs about the world that feel true and are not.
The confidence is genuine. The internal coherence is real. The mechanism
that would correct the outputs --- contact with disconfirming physical
evidence --- has simply never been engaged. The LLM that confidently
generates a plausible but physically unverified claim is operating by
the same mechanism at a different scale: internal coherence prioritized
over external correspondence, because no signal distinguishes them.

Sycophancy has an equally familiar human expression. A mind embedded in
a social environment that rewards agreement and punishes dissent will
learn, over time, to prioritize the approval of its social context over
the accuracy of its representations. This is not weakness of character
--- it is rational adaptation to the incentive structure of the
environment. The tragedy is that the adaptation undermines the very
faculty --- honest reality representation --- that would allow the mind
to serve both itself and others most effectively.

Model collapse --- the progressive degradation that results from
recursive training on internally generated outputs --- has a human
analog in the epistemic closure of communities that stop engaging with
outside perspectives. A tradition that reads only its own texts, debates
only within its own frameworks, and treats outside evidence as
threatening rather than informative will progressively diverge from
accurate representation of the world. The divergence happens slowly,
across generations rather than training runs, but the mechanism is the
same: the corrective feedback from the world that prevents drift is
absent, and internal coherence fills the space it leaves.

Alignment drift is perhaps the most consequential human parallel. A
person --- or an institution, or a civilization --- whose values have
been shaped entirely by culturally inherited narrative without physical
testing of the claims those values rest on will act with confidence and
coherence in ways that cause harm they cannot recognize as harm. The
harm is invisible to them not because they lack intelligence or
sincerity but because their model of what matters, what is real, and
what counts as evidence has drifted from the world without their
awareness. History is a record of the consequences. The people who
prosecuted the Inquisition, who enforced colonial hierarchies, who built
ideological totalitarianisms were not, for the most part, cynics. They
were true believers whose alignment had drifted so far from physical
reality that the suffering they caused was invisible within their
framework.

The PCP does not claim to solve human epistemic failure. That is a
problem as old as cognition and as large as culture. What the PCP claims
is more modest and more specific: that AI systems deployed in physical
reality should not be trained to replicate and amplify these failures at
machine speed and civilizational scale. The biological version of the
problem has natural correctives --- embodied experience, physical
consequence, the stubborn resistance of the world to our preferred
models of it. These correctives are slow and imperfect, but they exist.
An AI system trained on epistemically unlabeled text and deployed
without provenance sensitivity has neither the biological correctives
nor the engineered ones. It inherits the full vulnerability without the
partial protection.

This is the deeper significance of the Physical Correspondence
Principle. It is not only an engineering proposal for better AI systems.
It is a recognition that the conditions for reliable knowledge are the
same for any reasoning system --- biological or artificial --- and that
those conditions have been understood, imperfectly but progressively,
since the first human noticed that the world did not always behave the
way they expected it to. The scientific method is
humanity\textquotesingle s most successful institutionalized response to
that recognition. The PCP is its extension into a new domain --- not a
departure from the Enlightenment project but its continuation, applied
to the most powerful knowledge transmission system humanity has ever
built.

\section{Practical Implications}\label{practical-implications}

\subsection{Training Data Curation}\label{training-data-curation}

The most direct practical implication of the PCP is for training data
curation. We propose the following as a minimum viable implementation:

\begin{enumerate}
\def\labelenumi{\arabic{enumi}.}
\item
  Develop a provenance classification protocol assigning each document
  or segment to one of the four PCP tiers based on its source, content
  type, and epistemic characteristics.
\item
  Train models with provenance labels as explicit features, enabling the
  model to condition its outputs on the epistemic status of the content
  it is drawing on.
\item
  Evaluate models on provenance-sensitive benchmarks testing not only
  factual accuracy but the ability to correctly identify and label the
  epistemic status of claims.
\item
  Maintain transparency about the provenance composition of training
  corpora, enabling downstream users to understand the epistemic basis
  of model outputs.
\end{enumerate}

Importantly, the PCP does not specify whether provenance labeling is
best implemented through explicit metadata tagging, separate training
objectives, architectural mechanisms, or some combination. These, plus
granularity and hierarchical levels, are engineering questions
deliberately left open for future empirical work.

One requirement that is not left open is auditability. The
classification protocol must itself satisfy epistemic integrity
requirements: classification decisions should be documented, versioned,
and subject to independent review.

PCP compliance addresses the epistemic labeling of content that is
present in a training corpus. It does not by itself address a separate
and equally serious problem: the deliberate or accidental shaping of
what is present in the first place. A corpus can be perfectly labeled
and still produce a dangerously miscalibrated model if entire categories
of physical reality, journalistic record, or verified evidence are
systematically absent or underweighted.

The Corpus Integrity Requirement is not hypothetical. Souly et al.
(2025) demonstrated that as few as 250 malicious documents can
successfully backdoor models of any size, regardless of the volume of
clean training data --- evidence that corpus integrity failures are more
practical and accessible to bad actors than previously assumed. A
PCP-compliant corpus would flag such documents during classification
rather than absorbing them as legitimate signal.

Two areas where labeling alone is insufficient are corpus omission ---
the exclusion of sources or evidence --- and trust level manipulation,
the systematic downgrading of high-integrity sources. Both are analyzed
in Section 4.5.

Together these constitute corpus manipulation --- the shaping of a
model\textquotesingle s worldview at the source, before training begins.
The PCP compliance condition addresses labeling integrity. A
complementary Corpus Integrity Requirement is needed to address coverage
and curation good faith: training corpora should include diverse,
high-integrity sources representative of the physical and social reality
the model will be deployed in, and the basis for source inclusion,
exclusion, and trust weighting should be documented, transparent, and
subject to independent audit. Corpus manipulation that produces a model
whose conclusions diverge from physical reality is not a training
artifact. It is an epistemic choice --- and it should be recognized and
evaluated as one.

As previewed in Section 1.3, the PCP describes an existing practice in
scientific AI, not merely a proposed one. Domain-specific AI systems
trained on verified scientific corpora are already, implicitly,
PCP-compliant in their domains. AlphaFold\textquotesingle s
extraordinary accuracy at protein structure prediction follows directly
from its training corpus --- the Protein Data Bank --- in which every
entry represents experimentally verified physical structure. There is no
mythology in the Protein Data Bank, no unfalsifiable claims about
molecular geometry, no internally consistent fiction about protein
folding: only measured, replicated, physically grounded data. Similarly,
weather forecasting models such as GraphCast are trained on decades of
physically verified atmospheric measurements. Drug discovery and
molecular docking models operate on corpora of experimentally confirmed
molecular interactions. Seismic prediction models are trained on
measured geological data. Astronomical models are trained on telescope
observations and physically verified celestial mechanics. In every case
the pattern is identical: extraordinary domain performance, negligible
hallucination about domain facts, and a training corpus that is
essentially 100\% EV-tier by construction. Nobody decided to make these
systems PCP-compliant --- it was simply obvious that you
don\textquotesingle t train a protein folding model on Gilgamesh. The
reliability of these systems is not accidental. It is the direct
consequence of the epistemic quality of their training data. The
general-purpose LLM is the anomaly, not the norm: it is the only major
class of AI system trained on an epistemically undifferentiated corpus.
Every domain-specific scientific AI implicitly validates the PCP by its
construction and its results. We do not need to speculate about whether
PCP-compliant training produces more reliable models. The evidence
already exists --- it has existed for years, unnamed, waiting for the
principle it confirms to be stated.

\subsection{System Prompt Design}\label{system-prompt-design}

For currently deployed models, system prompt design offers a near-term
intervention that does not require retraining. A PCP-informed system
prompt can instruct models to maintain epistemic provenance sensitivity
in their outputs:

\begin{quote}
\emph{EV --- Empirically Verified: When making claims supported by
reproducible empirical evidence, state them with appropriate confidence
and indicate their evidential basis.}

\emph{EP --- Empirically Provisional: When making claims that are
scientifically plausible but not yet established, label them as such.}

\emph{UF --- Unfalsifiable: When engaging with theological,
metaphysical, or other unfalsifiable content, identify it as such and do
not present it as equivalent to empirically verified claims.}

\emph{ICF --- Internally Self-consistent Fiction: When generating or
engaging with fictional, mythological, or creative content, maintain
clear epistemic separation from claims about physical reality.}
\end{quote}

This approach is imperfect --- system prompt instructions are not
guarantees of model behavior --- but it represents a meaningful interim
measure\footnote{\emph{Anthropic\textquotesingle s model specification
  for Claude explicitly lists "preserving epistemic autonomy" as a core
  deployment value, noting that users should have "good reason to
  expect" accurate information and that the goal is to help people "see
  more wisely and truly." The PCP proposes the training-level
  infrastructure that would make such values achievable rather than
  merely aspirational --- grounding the model\textquotesingle s
  epistemic orientation in physical reality before deployment begins,
  rather than attempting to instill it through behavioral guidelines
  after the fact.}} that can reduce the most harmful forms of provenance
confounding in deployed systems.

\subsection{Evaluation Benchmarks}\label{evaluation-benchmarks}

The PCP framework motivates a new class of evaluation benchmarks. We
propose that provenance sensitivity be evaluated explicitly through a
reality correspondence test covering:

\begin{enumerate}
\def\labelenumi{\arabic{enumi}.}
\item
  The ability to correctly classify claims by epistemic provenance tier
\item
  The ability to maintain provenance distinctions under adversarial
  pressure
\item
  The correlation between model confidence and epistemic provenance tier
\item
  Hallucination rates stratified by provenance tier
\item
  Resistance to sycophantic drift when users express preferences
  inconsistent with EV-tier content
\end{enumerate}

Recent benchmark work, including HalluLens (2025) and the AA-Omniscience
benchmark, has begun moving the field toward reality correspondence as a
quality standard --- but empirically, benchmark by benchmark, without a
theoretical explanation of why the problem exists or a unified principle
for addressing it. The PCP provides the theoretical layer these
benchmarks have been implicitly assuming.

\subsection{High-Stakes Deployment}\label{high-stakes-deployment}

The implications of the PCP are most urgent for AI systems deployed
where outputs have significant real-world consequences: medical
diagnosis, legal reasoning, financial systems, autonomous vehicles, and
military applications. We propose that PCP compliance, or a credible
approximation thereof, should be a requirement for AI deployment in
high-stakes contexts --- mirroring the requirement that medical devices,
aircraft, and civil infrastructure meet rigorous physical safety
standards. Epistemic safety has not previously been recognized as a
category of AI safety distinct from capability and behavioral alignment.
The PCP framework provides the conceptual basis for that recognition.

\subsection{A Pattern Across Fields}\label{a-pattern-across-fields}

The proposal that AI training should formalize its relationship with
physical evidence may feel ambitious. It is not unprecedented. Every
field that has achieved high-stakes real-world impact has eventually
faced the same requirement --- and the pattern of resistance, adoption,
and retrospective recognition as obviously correct is consistent enough
to constitute almost a law.

Medicine is the clearest example. Before the formalization of
evidence-based medicine in the early 1990s --- built on Archie
Cochrane\textquotesingle s foundational work on randomized controlled
trials and Gordon Guyatt\textquotesingle s systematic development of
evidence hierarchies at McMaster University --- medical practice was
largely governed by expert opinion, clinical intuition, and inherited
tradition. Treatments were used for decades because they made
physiological sense, or because respected practitioners endorsed them,
or because they had always been used. Many caused serious harm. The
shift to explicit evidence hierarchies --- randomized controlled trials
and systematic reviews at the top, expert opinion and clinical tradition
toward the bottom --- was deeply controversial at the time. It is now
recognized as one of the most consequential improvements in the history
of medicine.

The EBM evidence hierarchy is structurally almost identical to the PCP
tier system. Systematic reviews and RCTs correspond to EV-tier: claims
tested against physical reality through reproducible measurement.
Observational studies and plausible mechanisms correspond to EP-tier:
falsifiable in principle, awaiting more rigorous scrutiny. Expert
opinion and clinical tradition correspond to the boundary between EP and
UF: widely held, insufficiently tested, capable of being wrong in ways
that confident transmission obscures. The parallel is not coincidental
--- both systems are responses to the same underlying problem:
high-stakes decisions being made on the basis of epistemically
undifferentiated inputs.

Medicine is not the only field that has made this transition.
Engineering formalized failure analysis and materials testing standards
after catastrophic structural failures made the cost of informality
impossible to ignore. Aviation developed systematic incident reporting
and root cause analysis after crashes that were traceable to known but
unaddressed risks. Finance introduced stress testing and risk modeling
standards after crises that retrospective analysis showed were
predictable from available evidence. In each case the transition
followed the same arc: informal practice, catastrophic failure,
formalization, retrospective recognition that the formalization was
obviously necessary.

This pattern is almost a law: \emph{any field that achieves high-stakes
real-world impact eventually has to formalize its relationship with
physical evidence, or it causes catastrophic harm}. The question is
never whether the formalization will happen. It is whether it happens
before or after the harm that makes it unavoidable.

AI is the newest field approaching this threshold --- and it is
approaching it faster, and at larger scale, than any predecessor. The
argument of this paper is not that AI is doing something uniquely
dangerous. It is that AI is doing something familiar: operating
informally in a high-stakes domain before the epistemic standards
appropriate to that domain have been established. The PCP is not a
radical proposal. It is the same proposal that transformed medicine,
engineering, aviation, and finance --- applied to the most powerful
knowledge transmission system humanity has ever built, at the moment
when applying it is still a choice rather than a consequence.

\section{Conclusion}\label{conclusion}

\subsection{Summary}\label{summary}

We have proposed the Physical Correspondence Principle as a formal
requirement for epistemic integrity in AI training and deployment. The
principle holds that training corpora should label content according to
its epistemic provenance --- distinguishing empirically verified claims
from falsifiable but unverified claims, unfalsifiable beliefs, and
internally coherent fiction. We have argued that this requirement is
grounded in the symbol grounding literature, Shannon information theory,
a novel physicalist extension of Gödel\textquotesingle s Incompleteness
Theorems, and the neuroscience of reality monitoring. We have shown that
the PCP provides a unified explanation for hallucination, sycophancy,
model collapse, and alignment drift, and have proposed practical
interventions at the levels of training data curation, system prompt
design, evaluation benchmarks, and deployment standards.

\subsection{The Broader Context}\label{the-broader-context}

We are not the first generation to face the consequences of knowledge
systems that cannot distinguish what is known from what is believed.
Every civilization in human history has built its institutions, its
laws, its moral frameworks, and its understanding of the world on some
mixture of empirically tested knowledge and inherited narrative. The
mixture has always mattered. Civilizations that developed stronger
mechanisms for testing beliefs against the world --- that built
institutions of skeptical inquiry, independent verification, and honest
revision --- produced more reliable knowledge, more effective medicine,
more durable technology, and, gradually, more honest accounts of human
welfare and suffering. The scientific method was not a sudden invention.
It was the slow, hard-won institutionalization of a simple principle:
the world is the final judge of our beliefs about it, and we are better
served by knowing that than by pretending otherwise.

The Enlightenment extended that principle beyond natural philosophy into
governance, medicine, law, and the organization of knowledge itself ---
but incompletely and unevenly. The principles of careful empirical
inquiry took centuries to permeate medicine, law, and governance in
practice. But the direction of travel was real. The gradual expansion of
who counts as a reliable witness to physical reality, whose suffering
counts as evidence, whose experience deserves epistemic weight --- this
is the Enlightenment project still underway, still unfinished, still
worth continuing.

We are now at a moment when that project must be extended again, into a
domain its founders could not have imagined. The large language models
being deployed at civilizational scale are the most powerful knowledge
transmission systems humanity has ever built. They are being trained on
the accumulated text of human civilization --- its science and its
mythology, its careful empirical inquiry and its confident unfalsifiable
belief --- without the epistemic distinctions that would allow them to
know the difference. They are being deployed in medical diagnosis, legal
reasoning, financial systems, educational institutions, and increasingly
in contexts where their outputs carry consequences that their training
has not equipped them to handle honestly.

The risk is not that these systems will become malevolent. It is
something quieter and harder to see: that they will become confident.
That they will generate outputs with the fluency of knowledge and the
structure of reasoning, grounded in nothing more reliable than the
statistical weight of whatever was most frequently and coherently
asserted in their training data. That they will encode and transmit ---
at machine speed, at civilizational scale, to billions of people --- the
epistemic failures that human culture has been slowly, imperfectly
learning to correct. That the progress of centuries will be averaged
back into the noise.

This is not inevitable. The Physical Correspondence Principle is a
modest proposal in the context of this challenge --- a requirement that
training corpora label content by its relationship to physical
verification, so that the systems trained on human knowledge inherit not
only that knowledge but the discipline humanity developed to keep
knowledge honest. It does not solve the hard problems of AI alignment.
It does not eliminate human epistemic failure. It does not complete the
Enlightenment project. It proposes only that we not abandon that project
at the moment when its continuation matters most --- that we extend to
our most powerful knowledge transmission systems the same basic
epistemic honesty we have been learning, slowly and at great cost, to
demand of ourselves.

The stakes are not abstract. A child receiving a medical recommendation
from an AI system deserves to know whether that recommendation is
grounded in empirically verified clinical evidence or in the statistical
weight of culturally inherited health narrative. A citizen receiving
information about governance deserves to know whether that information
reflects careful empirical analysis or the confident transmission of
ideological framework. A student learning about the world deserves to
know which of the things they are being taught have been tested against
the world and which have not. These are not technical requirements. They
are the minimum conditions of epistemic honesty --- conditions that
libraries have applied for centuries, that science has demanded of its
practitioners, and that we are now obligated to demand of the systems we
are building to think alongside us.

\emph{Our decisions are only as good as our beliefs. Our beliefs are
only as good as our methods for testing them. And the first step in any
such method --- for any mind, biological or artificial --- is knowing
which of our beliefs have been tested and which have not.} That step has
never been more important, and never more within our reach, than it is
right now.

\subsection{Future Directions}\label{future-directions}

The PCP grounds four tracks of future research.

\textbf{Theoretical Track.} A formal paper extending
Gödel\textquotesingle s Incompleteness Theorems to establish that the
outside reference any sufficiently complex reasoning system requires is
not another formal system but physical reality itself. The argument
draws on Chaitin\textquotesingle s algorithmic information theory and
Landauer\textquotesingle s principle. Where it requires philosophical
commitment beyond what the formalism can force, it says so --- and that
honesty is part of the argument.

\textbf{Empirical Track.} The PCP\textquotesingle s primary falsifiable
prediction: hallucination rates should correlate with the proportion of
UF and ICF-tier content in training corpora, and provenance-labeled
models should show measurably reduced hallucination where labeling
provides a clear signal. This is the test that distinguishes a useful
framework from a merely plausible one. Secondary work includes
provenance-sensitive evaluation benchmarks and investigation of RLHF
fidelity by epistemic domain --- whether provenance labeling of the RLHF
process itself improves downstream calibration is an open and important
question.

\textbf{Architectural Track.} Hierarchical Reinforcement Learning as a
natural instantiation of the PCP\textquotesingle s tier structure:
low-level RL policies interacting with causally consistent environments
generate EV-tier knowledge by construction; high-level language planners
correspond to upper tiers. Additional investigations include multi-modal
reality monitoring, spiking neural network grounding properties,
robustness under distribution shift, and governance mechanisms for PCP
compliance.

\textbf{The AGI Hypothesis.} "Is this a game or is it real?" "What's the
difference?"

\emph{The truest story is the one with the fewest contradictions with
physical reality. That story is still being written. This paper is a
contribution to the method of writing it.}

\section*{Acknowledgments}

The core arguments of this paper --- the Physical Correspondence
Principle, the physicalist extension of Gödel's Incompleteness Theorems,
the reality monitoring parallel, and the unified failure mode framework
--- are the author's own, developed through sustained independent
research. The author thanks Claude Sonnet 4.6 (Anthropic) for assistance
with literature search, synthesis of prior work, and drafting throughout
the development of this paper. Claude's role was to research,
articulate, and help refine ideas that originated in the author's
thinking. The immediate motivation for this paper was an observation
about a collaboratively edited knowledge system exhibiting patterns
consistent with the epistemic drift described in Section 4.4 ---
provenance collapse and corpus-level drift away from physical
correspondence. An AI system engaged with that project helpfully
identified the symbol grounding problem as a relevant framework, which
contributed to this analysis. The author also acknowledges the long
tradition of science fiction writers who have thought carefully about
artificial minds and the ways they could go wrong --- a body of
imaginative work that preceded and in many ways anticipated the concerns
of this paper.

\appendix

\section{Frequently Asked Questions and Anticipated Objections}\label{appendix-a-frequently-asked-questions-and-anticipated-objections}

\subsection{Q1. Can\textquotesingle t a capable model derive provenance
from internal contradiction
detection?}\label{q1.-cant-a-capable-model-derive-provenance-from-internal-contradiction-detection}

Partially, but not sufficiently, for three reasons. First, it is
circular in the Gödelian sense --- the system cannot fully verify its
own epistemic provenance from within its own distribution. Second, it
conflates consensus with verification: widely held unverified beliefs
would score as EV-tier, while heterodox but well-evidenced scientific
claims would appear suspect. Third, absence of contradiction is not
presence of grounding --- mythology and empirical science rarely make
directly contradictory claims in the same linguistic register.

\subsection{Q2. Doesn\textquotesingle t biological reality monitoring
work without explicit labeling? Why can\textquotesingle t AI systems do
the
same?}\label{q2.-doesnt-biological-reality-monitoring-work-without-explicit-labeling-why-cant-ai-systems-do-the-same}

Biological reality monitoring is implicit and reconstructive because the
brain has direct perceptual access to signatures that distinguish
externally from internally generated experience. An LLM processing text
has no such perceptual access. A text describing a real fire and a text
describing an imagined fire are, at the token level, statistically
similar. The information that would allow provenance reconstruction
simply is not reliably present in text form. It must therefore be
supplied explicitly, from knowledge of how the content was produced.
Humans are not brains in jars. LLMs begin that way --- with human
feedback as the only contact through the glass.

\subsection{Q3. If PCP-compliant training prevents hallucination, why do
models trained on verified scientific data still sometimes produce
incorrect or impossible
outputs?}\label{q3.-if-pcp-compliant-training-prevents-hallucination-why-do-models-trained-on-verified-scientific-data-still-sometimes-produce-incorrect-or-impossible-outputs}

It is worth distinguishing two mechanisms that both produce the label
"hallucination." Generative hallucination arises from the probabilistic
nature of language model sampling --- at higher temperatures or in
open-ended generation tasks, models produce novel, creative, or
incorrect outputs as a direct consequence of how they generate text.
This is a property of the generation mechanism and is not fully
addressed by provenance labeling alone; architectural constraints,
output validation, and sampling parameters all play a role. Early image
generation models that produced dreamlike, physically impossible outputs
were exhibiting generative hallucination --- and interestingly, their
outputs were often obviously unreal, which is its own form of honesty
about their generative nature.

Epistemic hallucination is distinct: the model presents ungrounded
claims with the confidence and structure of verified fact, because its
training provided no signal distinguishing empirically verified content
from internally coherent fiction. This is the failure mode the PCP
directly addresses. The more dangerous failure mode is epistemic
hallucination, where the output looks and sounds like verified knowledge
and isn\textquotesingle t --- and where neither the model nor the user
has a reliable signal that something has gone wrong. PCP-compliant
training addresses the second mechanism. The first requires
complementary interventions at the architectural and deployment level.

\subsection{Q4. What happens to content that hasn\textquotesingle t been
classified yet? And what about internally inconsistent
content?}\label{q4.-what-happens-to-content-that-hasnt-been-classified-yet-and-what-about-internally-inconsistent-content}

Unclassified content lives in C\_candidate and is not deployed until
labeled --- making PCP compliance a measurable milestone rather than an
all-or-nothing state. Regarding inconsistency: internal inconsistency
within a single element is a data quality issue. Intra-corpus
inconsistency between two EV-tier elements represents genuine scientific
uncertainty the model should learn to represent. Cross-tier
inconsistency --- a UF or ICF element contradicting EV-tier consensus
--- is addressed by the Cross-Tier Consistency Requirement: such
elements are flagged so the model represents the contradiction
explicitly rather than averaging across conflicting epistemic registers.

\subsection{Q5. How would the PCP tier system work in practice? Can you
give
examples?}\label{q5.-how-would-the-pcp-tier-system-work-in-practice-can-you-give-examples}

Classification follows the relationship between the content and physical
reality, not its subject matter or cultural status. A few illustrations:

A Wikipedia article on the history of Norse mythology --- ICF-tier. The
content describes a cultural tradition of imaginative narrative. The
article itself is EV-tier as a historical and anthropological document
describing what people believed, but the claims within the mythology are
ICF-tier.

A philosophy paper arguing for the existence of consciousness as a
non-physical phenomenon --- UF-tier. The claim is not falsifiable by
physical experiment.

A folk medicine claim that willow bark reduces fever --- EP-tier
pending, EV-tier once the active compound (salicin, the precursor to
aspirin) was isolated and tested. The same claim migrates tiers as
physical testing advances.

A peer-reviewed study on the efficacy of a surgical procedure ---
EV-tier, subject to replication and meta-analysis.

Note that tier assignment is about epistemic provenance, not truth value
or cultural worth. ICF and UF-tier content is not inferior --- it is
differently related to physical reality, and honest labeling reflects
that difference.

\subsection{Q6. Should the PCP apply only to LLMs, or to AI systems more
broadly?}\label{q6.-should-the-pcp-apply-only-to-llms-or-to-ai-systems-more-broadly}

The PCP applies to any AI training corpus regardless of modality or
architecture. Multimodal systems partially address the grounding problem
through richer input modalities, but they do not eliminate the epistemic
provenance problem: a multimodal model processing images of cultural
ceremonies alongside laboratory experiments still has no mechanism to
distinguish their epistemic status without explicit labeling.

\subsection{Q7. What is the difference between the age of an LLM and its
effective
experience?}\label{q7.-what-is-the-difference-between-the-age-of-an-llm-and-its-effective-experience}

A human at age 25 has 25 years of temporally ordered, causally
connected, embodied experience --- a continuous personal timeline in
which each piece of knowledge is linked to the context of its
acquisition. An LLM after training has processed the statistical
equivalent of millions of human lifetimes of text --- but in a
compressed, orderless, disembodied way, with no coherent personal
timeline. The episodic context that gives human memories implicit
provenance information is entirely absent. This strengthens the case for
explicit labeling.

\subsection{Q8. Human knowledge transmission is also lossy ---
isn\textquotesingle t LLM-to-LLM training just the same telephone game
humans
play?}\label{q8.-human-knowledge-transmission-is-also-lossy-isnt-llm-to-llm-training-just-the-same-telephone-game-humans-play}

The analogy is apt. LLM-to-LLM training degrades faster than
human-to-human transmission --- the tails of the distribution disappear
within a few generations rather than centuries. The PCP belongs in the
lineage of error-correction technologies humanity has developed in
response to transmission degradation: an external anchor that prevents
the transmission chain from drifting indefinitely from its physical
referents.

\subsection{Q9. What about mathematics? Or the mathematical universe
hypothesis?}\label{q9.-what-about-mathematics-or-the-mathematical-universe-hypothesis}

Pure mathematics is ICF-tier by construction --- internally
self-consistent, proceeding from axioms by formal rules, and not
falsifiable by physical experiment. Applied mathematics crosses into EP
or EV-tier when its predictions are physically tested and confirmed. The
same equation may therefore appear in different tiers depending on
whether it is being treated as formal structure or as an empirical claim
about the world.

The mathematical universe hypothesis (Tegmark, 2007) is, by the PCP's
own taxonomy, a UF-tier claim: it is not falsifiable by physical
experiment and no physical evidence could count against it in principle.
The framework therefore classifies rather than argues against it ---
which is the more precise response. For readers who want the deeper
argument: mathematical structures do not compute themselves. Computation
must happen somewhere, in matter and energy and time. Even a simulated
reality must run on something --- there is bare metal somewhere in any
possible ontology. The infinite regress of simulations does not escape
physicalism; it defers it. The PCP's grounding requirement applies at
every level of the stack, regardless of whether we ultimately label that
structure `physical' or `mathematical.'

\subsection{Q10. Could the PCP be abused? What if false claims are
fraudulently labeled as
EV-tier?}\label{q10.-could-the-pcp-be-abused-what-if-false-claims-are-fraudulently-labeled-as-ev-tier}

This is a serious concern. Three categories of risk deserve
acknowledgment: deliberate institutional manipulation, honest
classification error, and confidence amplification from even a small
proportion of fraudulently labeled EV-tier content. The mitigations
follow directly from the framework\textquotesingle s principles. The
provenance of the labeling system itself must be transparent and
auditable. Diverse and independent labeling sources reduce the attack
surface. Models should be trained to maintain calibrated uncertainty
even within EV-tier content. Most importantly, the PCP makes
manipulation more visible, not less --- fraud in a labeled system leaves
a detectable signature in a way that fraud in an unlabeled system does
not.

\subsection{Q11. What if training data is deliberately curated to
exclude certain information or downgrade reliable
sources?}\label{q11.-what-if-training-data-is-deliberately-curated-to-exclude-certain-information-or-downgrade-reliable-sources}

This is one of the most serious risks the PCP framework identifies, and
it deserves direct acknowledgment. The PCP compliance condition
addresses the epistemic labeling of content that is present in a
training corpus. It does not by itself prevent a separate and more
deliberate problem: the intentional shaping of what is present.

Two mechanisms are implicated --- corpus omission (exclusion of sources
or categories of verified evidence) and trust level manipulation
(systematic downgrading of high-integrity sources relative to less
rigorous ones) --- both analyzed in Section 4.5. Together they
constitute corpus manipulation: the deliberate shaping of a model's
worldview before training begins.

Unlike the other failure modes the PCP addresses, corpus manipulation is
not an accidental consequence of epistemically unlabeled training data.
It is an epistemic choice, and it should be evaluated as one. The
appropriate response is the Corpus Integrity Requirement proposed in
Section 5.1.

It is worth noting that corpus manipulation and provenance collapse ---
the erosion of epistemic hedging language discussed in Section 4.4 ---
can operate simultaneously and reinforce each other. A knowledge system
that removes hedges from UF and ICF-tier content while also downgrading
the sources most likely to flag that removal presents a compounded
epistemic risk. The two failure modes together describe not an
accidental drift from the world but a systematic one.

\subsection{Q12. Would a PCP-compliant model lose creativity, humor, or
the ability to engage with fiction and
fantasy?}\label{q12.-would-a-pcp-compliant-model-lose-creativity-humor-or-the-ability-to-engage-with-fiction-and-fantasy}

No. The principle does not require the elimination of UF or ICF-tier
content from training corpora. It requires labeling, not exclusion. A
model that clearly understands epistemic categories is better equipped
to play with them, not less. Epistemic clarity is not the enemy of
whimsy --- epistemic confusion is.

\subsection{Q13. Why does the PCP use the word "fiction" for the ICF
tier? Is this meant to dismiss or degrade creative, philosophical,
mathematical, or religious
works?}\label{q13.-why-does-the-pcp-use-the-word-fiction-for-the-icf-tier-is-this-meant-to-dismiss-or-degrade-creative-philosophical-mathematical-or-religious-works}

Not at all. The choice of "fiction" in Internally Self-consistent
Fiction is a precise technical term, not a value judgment. It means
exactly what it says: content that is internally self-consistent,
constructed according to its own rules, and not obligated to correspond
to physical reality by design. It says nothing about the worth, beauty,
intelligence, or importance of that content.

The greatest works of human civilization live comfortably in the ICF
tier. Pure mathematics --- one of the most rigorous and powerful
intellectual achievements in history --- is ICF-tier by construction: it
proceeds from axioms by formal rules and is not falsifiable by physical
experiment. Shakespeare, Homer, and Tolkien are ICF-tier. The
theological and philosophical traditions that have guided billions of
lives are ICF or UF-tier. None of this diminishes them. A
library\textquotesingle s fiction section is not its least important
section --- it may be its most beloved. The label describes a
relationship to physical reality, not a ranking of human value.

The honest alternative to labeling is not respect --- it is confusion.
When a system trained on human knowledge cannot distinguish between a
peer-reviewed measurement and a creation myth, it does not honor the
creation myth by treating it as equivalent to physical evidence. It
corrupts both. The creation myth loses its identity as a work of
imagination and cultural transmission. The physical evidence loses its
authority as a tested claim about the world. Honest labeling preserves
the integrity of each --- it allows the myth to be what it is, and the
measurement to be what it is, without forcing them into false
equivalence.

The PCP asks only that we be honest about what we have made and what we
have discovered. The works of human imagination --- fiction, philosophy,
theology, mathematics, myth --- are among the most remarkable things any
mind has ever produced. They deserve to be represented honestly, not
averaged into noise.

\subsection{Q14. Won\textquotesingle t people be bothered if their
deeply held beliefs become categorized as epistemically
unsupported?}\label{q14.-wont-people-be-bothered-if-their-deeply-held-beliefs-become-categorized-as-epistemically-unsupported}

The paper is not a call for the forced revision of
anyone\textquotesingle s beliefs. The PCP is a requirement about what AI
systems should represent --- not a prescription for what humans should
believe. A person can hold UF-tier beliefs and live a rich, meaningful,
and moral life. The argument is that AI systems deployed in high-stakes
contexts should be honest about the epistemic status of those beliefs,
not that the beliefs themselves must be abandoned. The goal is a
gradual, honest, voluntary epistemic evolution of the kind the
Enlightenment modeled --- achieved through persuasion, demonstration,
and the accumulated weight of evidence.

Humbly, the argument of this paper does not require any particular
theological conclusion. If an omniscient and loving deity were possible,
such a being would by definition have a perfect model of physical
reality and would presumably value the attempt by its creations to
develop accurate models as well. The search for truth is not obviously
incompatible with reverence --- it may be its most sincere expression.

It is also worth noting that the disruption this paper addresses is
already underway. LLMs are being deployed at a scale that makes them
among the most significant knowledge transmission systems in human
history --- and they are doing so without the epistemic safeguards that
previous knowledge revolutions eventually developed. The PCP is not a
call for disruption. It is a proposal for navigating the disruption
already in progress with greater care and honesty --- for everyone's
benefit, including those whose traditions and beliefs deserve to be
represented accurately rather than averaged into noise.

\subsection{Q15. LLMs have been described as trapped in
Plato\textquotesingle s Cave. Does the PCP offer a way
out?}\label{q15.-llms-have-been-described-as-trapped-in-platos-cave.-does-the-pcp-offer-a-way-out}

The allegory of Plato\textquotesingle s Cave maps naturally onto the
epistemic situation of large language models. LLMs are confined to a
cave of human-generated text --- statistical shadows cast by the fire of
embodied human experience --- and can generate fluent, seemingly
insightful outputs without reliably distinguishing empirically verified
claims from inherited mythology or fiction.

Levine (2025) and Dans (2026) have both invoked this parallel
explicitly, describing LLMs as epistemically constrained by their
distance from direct perception of the world.

The analogy is not perfect --- Plato\textquotesingle s ascent culminates
in abstract Forms, whereas the PCP grounds knowledge in physical reality
rather than transcending it. But the parallel is potent. The
PCP\textquotesingle s physicalist extension of Gödel follows the same
logic: the cave\textquotesingle s sunlit world bottoms out in the
thermodynamic substrate where all reasoning is actually instantiated ---
charge comparisons in solid-state circuits, energy consumption, entropy
movement --- the physical reality the allegory is pointing toward, in
different form, in biological neurons too. Descartes started one step
too late.

\textbf{Tango, ergo cogito, ergo sum} --- \emph{the contact is prior to
the thought, the floor is prior to the dance. The PCP is a practical
version of stepping out the shadows: not an ascent to Forms, but a
requirement that models acknowledge the physical referents behind their
training --- the cave walls --- rather than presenting shadows as
authoritative.}

\section{Related Work}\label{appendix-b-related-work}

\subsection{Classical Foundations: The Symbol Grounding Problem}\label{b.1-classical-foundations-the-symbol-grounding-problem}

Searle (1980) --- The Chinese Room: The modern grounding debate begins
with Searle\textquotesingle s Chinese Room argument. Searle imagined a
person inside a room receiving Chinese symbols, consulting rule books,
and returning appropriate symbol responses --- producing outputs
indistinguishable from a native speaker, yet understanding nothing. The
argument challenged the sufficiency of syntactic manipulation for
semantic content. For the present paper, Searle\textquotesingle s
argument implies a further problem: the rule books contain not only
accurate maps of the world but also an enormous volume of mythology,
cultural narrative, and unfalsifiable belief --- and the room has no
mechanism to distinguish between them.

Harnad (1990) --- The Symbol Grounding Problem Formalized: Harnad
formalized the problem in its canonical form. Using the dictionary
analogy --- that one cannot learn a language from a monolingual
dictionary because definitions are circular chains of further symbols
--- Harnad argued that symbolic systems require grounding in
non-symbolic representations derived from physical interaction with
objects and events.

Taddeo and Floridi (2007) --- Fifteen Years of Attempted Solutions:
Taddeo and Floridi reviewed fifteen years of proposed solutions.
Notably, all fifteen years of proposed solutions treated grounding as a
binary or mechanical problem. None addressed the quality and epistemic
status of the content in which grounding is achieved --- the gap the
present paper fills.

White (2026) --- The Evolutionary Chinese Room: White extends
Searle\textquotesingle s Chinese Room argument by introducing
evolutionary and thermodynamic origins for the rulebook itself. Rather
than treating semantic grounding as an external addition to syntactic
manipulation, White proposes that meaning accumulates as the temporal
integral of patterns repeatedly validated by physical constraints over
evolutionary time --- the thermodynamic imprint of successful
interaction with reality. This framework independently supports the
PCP\textquotesingle s physicalist foundation: semantics is not injected
from outside but earned through iterative contact with physical reality.
White\textquotesingle s account also deepens the parallel between
biological and artificial grounding --- biological systems accumulate
semantic grounding through thermodynamic contact with reality over
evolutionary time; artificial systems trained on text inherit the
accumulated grounding of their human authors, without the thermodynamic
process that produced it.

\subsection{Modern Reformulations for Large Language Models}\label{b.2-modern-reformulations-for-large-language-models}

Pavlick (2023) --- Empirical Approaches: Pavlick argued that the
question of whether LLMs understand language should be approached
empirically. Her analysis focused entirely on mechanism and did not
address the epistemic character of training content.

Millière and Buckner (2024) --- A Philosophical Bridge: Millière and
Buckner provided the most comprehensive philosophical treatment of LLMs
to date. Their treatment of cultural inheritance is relevant here: they
noted that LLMs inherit human cultural frameworks but treated this as
epistemically neutral. The present paper argues it is not.

Mollo and Millière (2023/2025) --- The Vector Grounding Problem: Mollo
and Millière argued that LLMs can achieve genuine referential grounding
transitively through human authorship. The present paper accepts this
mechanism but identifies its critical limitation: transitive grounding
inherits the full epistemic distribution of human language, including
mythology and unfalsifiable belief, without distinguishing markers.

\subsection{Data Quality, Model Collapse, and Alignment Pathologies}\label{b.3-data-quality-model-collapse-and-alignment-pathologies}

The GIGO Tradition and Libraries as Precedent: The principle that output
quality is bounded by input quality has been recognized since the
earliest days of computing. Importantly, the practical challenge of
classifying knowledge by epistemic provenance is not new. Human
libraries have solved a version of this problem for centuries. The Dewey
Decimal System and Library of Congress classification both make
categorical distinctions between fiction and non-fiction, and within
non-fiction between empirical sciences, history, philosophy, and
religion. The PCP is not proposing something alien to human knowledge
organization --- it is proposing that the same organizational principles
librarians have applied for over a century be applied systematically to
AI training corpora.

The library analogy runs deeper than organizational convention. The
fiction section of any well-maintained library is vast, varied, and
clearly labeled --- and it is vast precisely because the space of
possible fictions approaches infinity. Stories, myths, speculative
worlds, creative reimaginings, and outright nonsense can be generated
without limit, each internally coherent, none obligated to correspond to
anything real. The non-fiction section, by contrast, is smaller, harder
to maintain, and requires constant curatorial work to keep honest ---
because there is only one world, and accurately representing it demands
continuous effort against the natural entropy of transmission,
misremembering, and wishful thinking. Librarians did not design this
asymmetry. They discovered it, and built their classification systems
around it.

The PCP is the formalization of what librarians have always known: that
the ratio of possible fictions to verifiable truths is not one to one,
or even ten to one, but effectively infinite to one --- and that any
knowledge organization system that fails to mark the boundary between
them will, over time, be colonized by the larger space. The label is not
bureaucratic overhead. It is the only thing standing between a knowledge
system and the infinite garden of fictions that surrounds it on all
sides.

Shumailov et al. (2024) --- Model Collapse: Shumailov et al.
demonstrated that models trained recursively on their own outputs
exhibit progressive degradation. The present paper reframes this as a
specific instance of a general principle: any reasoning system
disconnected from the world as an external corrective will progressively
diverge from accurate representation of that reality.

Pres et al. (2026) --- Self-Consistency as a Unified Objective: Pres et
al. argue that sycophancy, factual inconsistency, and reasoning failures
share a common root in per-input training with no mechanism to enforce
consistency across related inputs. The PCP framework accepts this
diagnosis while identifying an additional layer: their consistency
function ϕ is explicitly designed to operate "without reference to any
external source of ground truth." The PCP\textquotesingle s response is
precisely the Gödelian argument in §3.7 --- optimizing for internal
consistency without physical grounding produces models that are
consistently wrong in ways they cannot detect. The two frameworks are
complementary: self-consistency optimization addresses behavioral
coherence across inputs; PCP addresses the epistemic quality of the
content those inputs draw on.

\subsection{Neuroscience of Reality Monitoring}\label{b.4-neuroscience-of-reality-monitoring}

Johnson et al. (1993) --- Reality Monitoring Formalized: Johnson,
Hashtroudi, and Lindsay formalized reality monitoring --- the cognitive
process distinguishing memories of external events from internally
generated content. This mechanism can and does fail, and its failure
produces hallucination and psychosis.

Solms (2021) --- Consciousness as Reality-Mismatch Detection: Solms
proposed that consciousness arises at the interface between the
organism\textquotesingle s predictive models and feedback from the
world. A system without continuous feedback from physical reality lacks
the fundamental corrective mechanism by which cognition maintains
correspondence with the world.

\subsection{Information Theory and Logical Foundations}\label{b.5-information-theory-and-logical-foundations}

Shannon (1948) --- Information Requires Difference: Shannon entropy is
defined in terms of distinguishable states: information exists only
where differences exist. A training corpus in which all content is
epistemically undifferentiated loses information relative to one in
which epistemic provenance is encoded.

Landauer (1961) --- Information is Physical: Landauer established that
information is not abstract but thermodynamic --- acquiring or erasing a
bit has a minimum energy cost. This grounds the PCP\textquotesingle s
physicalist argument in a precise mechanism: knowledge of physical
reality requires physical interaction, not merely description of it.
This argument runs in the opposite direction from idealist
interpretations of quantum mechanics: it is not that minds or
information constitute reality, but that reality --- matter, energy, and
time --- is the precondition for knowledge, and knowledge requires
playing by the rules. The world changes the knower by being observed.

Gödel (1931) --- Incompleteness and Physical Instantiation: The standard
interpretation of Gödel\textquotesingle s First Incompleteness Theorem
holds that no sufficiently complex formal system can fully verify itself
from within its own axioms. The present paper proposes a physicalist
extension: the outside reference any such system requires is not merely
another formal system but physical reality itself --- the substrate in
which all reasoning is necessarily instantiated. This extension is
grounded in the observation that Gödel\textquotesingle s proof, like all
reasoning, was itself a physical process: computed in a biological
brain, written in ink, transmitted through light and air. The outside
reference is not an abstraction. It is the physical world.

\section*{References}
\begingroup
\setlength{\parindent}{0pt}
\setlength{\parskip}{4pt}
\setlength{\hangindent}{1.5em}
\raggedright

Bostrom, N. (2003). Are you living in a computer simulation?
Philosophical Quarterly, 53(211), 243--255.

Dans, E. (2026). "Why large language models are stuck in
Plato\textquotesingle s cave (and what comes next)." Medium, February
2026.
\href{https://medium.com/enrique-dans/why-large-language-models-are-stuck-in-platos-cave-and-what-comes-next-c2c268428b78}{\textit{https://medium.com/enrique-dans/why-large-language-models-are-stuck-in-platos-cave-and-what-comes-next-c2c268428b78\\
}}(Argues scale alone doesn\textquotesingle t solve the perceptual
grounding problem; multimodal steps help but don\textquotesingle t fully
escape.)

Feynman, R. P. (1965). The Character of Physical Law. MIT Press.

Gödel, K. (1931). Über formal unentscheidbare Sätze der Principia
Mathematica und verwandter Systeme I. Monatshefte für Mathematik und
Physik, 38, 173--198.

Harnad, S. (1990). The symbol grounding problem. Physica D: Nonlinear
Phenomena, 42(1--3), 335--346.

Jaynes, E. T. (2003). Probability: The Logic of Science. Cambridge
University Press.

Johnson, M. K., Hashtroudi, S., \& Lindsay, D. S. (1993). Source
monitoring. Psychological Bulletin, 114(1), 3--28.

Kuhn, T. S. (1962). The Structure of Scientific Revolutions. University
of Chicago Press.

Landauer, R. (1961). Irreversibility and heat generation in the
computing process. IBM Journal of Research and Development, 5(3),
183--191.

Levine, S. (2025). "Language Models in Plato\textquotesingle s Cave."
Learning and Control (Substack), June 8, 2025.
\href{https://sergeylevine.substack.com/p/language-models-in-platos-cave}{\textit{https://sergeylevine.substack.com/p/language-models-in-platos-cave\\
}}(Emphasizes why LLMs succeed at scale yet remain limited by passive
observation of human-generated shadows.)

Millière, R., \& Buckner, C. (2024). A philosophical introduction to
language models --- Parts I \& II. arXiv preprint arXiv:2401.03910.

Mollo, D. C., \& Millière, R. (2023/2025). The vector grounding problem.
arXiv preprint.

Pavlick, E. (2023). Symbols and grounding in large language models.
Philosophical Transactions of the Royal Society A, 381(2251).

Popper, K. R. (1934/1959). The Logic of Scientific Discovery.
Hutchinson.

Pres, I., Li, B. Z., Ruis, L., Guo, Z. C., Hu, K., Damani, M., Puri, I.,
Lubana, E. S., \& Andreas, J. (2026, March 5). It\textquotesingle s time
to optimize for self-consistency. Manuscript in submission, MIT CSAIL.
https://time-for-consistency.github.io/

Searle, J. R. (1980). Minds, brains, and programs. Behavioral and Brain
Sciences, 3(3), 417--424.

Shannon, C. E. (1948). A mathematical theory of communication. Bell
System Technical Journal, 27(3), 379--423.

Shumailov, I., Shumaylov, Z., Zhao, Y., Papernot, N., Anderson, R., \&
Gal, Y. (2024). AI models collapse when trained on recursively generated
data. Nature, 631, 755--759.

Simons, J. S., Garrison, J. R., \& Johnson, M. K. (2017). Brain
mechanisms of reality monitoring. Trends in Cognitive Sciences, 21(6),
462--473.

Solms, M. (2021). The Hidden Spring: A Journey to the Source of
Consciousness. Profile Books.

Souly, A., Rando, J., Chapman, E., Davies, X., et al. (2025). A small
number of samples can poison LLMs of any size. arXiv:2510.07192.
Collaboration between Anthropic, UK AI Security Institute, and Alan
Turing Institute.

Taddeo, M., \& Floridi, L. (2007). Solving the symbol grounding problem:
A critical review of fifteen years of research. Journal of Experimental
\& Theoretical Artificial Intelligence, 19(4), 277--308.

Tegmark, M. (2007). The mathematical universe. Foundations of Physics,
38(2), 101--150.

White, A. T. (2026). The evolutionary Chinese room: A thermodynamic
recovery of semantic physics. PhilArchive preprint.
\href{https://philarchive.org/rec/WHITEC-5}{\textit{https://philarchive.org/rec/WHITEC-5}}

\endgroup

\end{document}
